\section{Хроматичность при синхротронном движении}
\label{sec:rf}

С ВЧ $\delta$ перестаёт быть параметром: она колеблется с синхротронной частотой, малой по сравнению с бетатронной. Бетатронная фаза набирается с мгновенной частотой $\nu(\delta(t))$, и наблюдаемое число колебаний есть её среднее по синхротронному периоду. Линейная хроматичность при усреднении выпадает, остаётся член второго порядка. Поэтому нормальная форма 6D (\texttt{TS} в COSY) даёт не $\xi$, а усреднённую вторую хроматичность. Линейную $\xi$ можно восстановить, сдвинув равновесную энергию --- неподвижную точку 6D-отображения --- смещением частоты ВЧ. Глава связывает это с равновесным сдвигом~\cite{EqNote}.

\subsection{Медленная переменная и усреднение}

Продольное движение (см.~\cite{EqNote}, глава~1) при $\nu_s\ll\nu_{x,y}$ даёт
\begin{equation}
\label{eq:delta-t}
\delta(t)=\delta_\eq+\delta_m\cos\psi_s+O(\delta_m^2),
\qquad
\langle\delta\rangle=\delta_\eq,
\qquad
\langle\delta^2\rangle=\delta_\eq^2+\frac{\delta_m^2}{2}.
\end{equation}
За один бетатронный период $\delta$ почти не меняется, и бетатронное движение адиабатически следует за замкнутой орбитой $x_{\mathrm{co}}(s;\delta(t))$. Бетатронная фаза за $N$ оборотов равна $2\pi\sum_k\nu(\delta_k)$, поэтому средняя частота --- это среднее мгновенной (метод усреднения~\cite{Bogoliubov1961}):
\begin{equation}
\label{eq:nu-avg}
\bar\nu=\bigl\langle\nu(\delta)\bigr\rangle
=\nu_0+\xi\,\delta_\eq+\frac{\nu^{(2)}}{2}\Bigl(\delta_\eq^2+\frac{\delta_m^2}{2}\Bigr)+\dots
\end{equation}
Модуляция частоты с синхротронной частотой даёт боковые полосы $\nu\pm k\nu_s$, но среднюю частоту не сдвигает: фаза есть интеграл частоты, и колеблющаяся часть $\xi\delta_m\cos\psi_s$ в среднее не входит. Для частицы около синхронной ($\delta_\eq\approx0$) линейная хроматичность из $\bar\nu$ выпадает; остаётся амплитудная зависимость от синхротронного действия через $\nu^{(2)}$.

\subsection{Что даёт нормальная форма 6D}

В 6D нормальной форме гамильтониан в переменных действие--угол зависит только от действий,
\begin{equation}
\label{eq:nf6}
\tilde{\mathcal{H}}=\mu_xJ_x+\mu_yJ_y+\mu_sJ_s
+\tfrac12\sum_{i,k}c_{ik}J_iJ_k+\dots,
\qquad
\nu_x(J)=\frac{1}{2\pi}\frac{\partial\tilde{\mathcal{H}}}{\partial J_x}.
\end{equation}
Синхротронное действие связано с амплитудой: $\langle\delta^2\rangle_{\psi_s}=\gamma_LJ_s$, где $\gamma_L$ --- продольная функция Твисса для $\delta$. Сравнение с \eqref{eq:nu-avg} при $\delta_\eq=0$ даёт
\begin{equation}
\label{eq:dnu-dJs}
\frac{\partial\nu_x}{\partial J_s}=\frac{\nu_x^{(2)}}{2}\,\gamma_L,
\end{equation}
т.\,е. перекрёстный амплитудный сдвиг <<бетатрон--синхротрон>>, который выдаёт \texttt{TS}, пропорционален второй хроматичности. Линейная $\xi$ в \eqref{eq:nf6} не появляется: в 6D нет параметра, по которому её можно было бы взять, а член $\xi\,\delta\,\mathcal{H}$ с колеблющимся $\delta$ нормальная форма убирает как нерезонансный. При $V\to0$ величина \eqref{eq:dnu-dJs} в нормировке COSY стремится к нулю вместе с $\nu_s$, а не к $\xi$. Если в 6D-постановке задать энергию ещё и параметром (\texttt{OV 3 3 1}), параметр в отображение не попадает: энергия уже есть координата, и производная по параметру равна нулю.

\subsection{Восстановление $\xi$ сдвигом неподвижной точки}

Сместим частоту ВЧ: $\omega_\rf\to\omega_\rf(1+\epsilon)$. Синхронной становится частица, у которой период обращения короче в $1+\epsilon$ раз, т.\,е. $\tau(\delta^*)=-\epsilon$, и в линейном порядке
\begin{equation}
\label{eq:dstar}
\delta^*=-\frac{\epsilon}{\eta_0}.
\end{equation}
Неподвижная точка 6D-отображения сдвигается в $(x,p_x)=(D_1,D_1')\delta^*$, $\delta=\delta^*$. Её якобиан $6\times6$ имеет три пары собственных значений $e^{\pm i\mu}$; поперечные частоты $\nu_{x,y}(\delta^*)$ считаются на смещённой орбите, и
\begin{equation}
\label{eq:xi-fp}
\xi=\frac{\mathrm{d}\nu_{x,y}}{\mathrm{d}\delta^*}\Big|_{\delta^*=0}.
\end{equation}
Это то же определение \eqref{eq:xi-trace}, только неподвижная точка находится в 6D, а $\delta$ задаётся частотой ВЧ. Синхротронная связь через дисперсию меняет собственные поперечные частоты якобиана лишь в высших порядках по малой связи, поэтому \eqref{eq:xi-fp} от напряжения практически не зависит и при $V\to0$ переходит в 4D-результат \texttt{OV 3 2 1}. Тот же приём --- измерение $\xi$ по сдвигу частот при изменении частоты ВЧ --- используется на действующих кольцах.

\subsection{Связь с равновесным сдвигом}

В~\cite{EqNote} показано, что усреднённое по бетатронной фазе удлинение пути равно
\begin{equation}
\label{eq:Lambda}
\Lambda=-\frac{2\pi}{C_s}\bigl(J_x\xi_x+J_y\xi_y\bigr),
\qquad
\delta_\eq=-\frac{1}{\eta_0}\Bigl(\frac{\eta_1\delta_m^2}{2}+\Lambda\Bigr).
\end{equation}
Хроматичность входит сюда как полная, со всеми элементами и краями: вывод~\cite{EqNote} опирается только на гамильтонову форму и статичность поля. Поэтому для $\delta_\eq$ нужно брать $\xi_y$ с членом F \eqref{eq:xiF}; модель <<только тело>> даёт в вертикальной плоскости слишком отрицательную $\xi_y$ и, следовательно, завышенное $\Lambda$. С другой стороны, \eqref{eq:nu-avg} показывает обратную связь: сдвиг $\delta_\eq$ \eqref{eq:Lambda} даёт сдвиг средней частоты $\xi\,\delta_\eq$, т.\,е. член второго порядка по амплитудам, того же порядка, что $\nu^{(2)}\delta_m^2/4$.
